\begin{lemma}[formulaLtOne]
\uses{VerticalIntegral}
\label{formulaLtOne}
\lean{Perron.formulaLtOne}
% at /home/runner/work/trough/trough/library/PrimeNumberTheoremAnd/PerronFormula.lean:731.0-772.58
\leanok
For $x>0$, $\sigma>0$, and $x<1$, we have
  $$
  \frac1{2\pi i}
  \int_{(\sigma)}\frac{x^s}{s(s+1)}ds =0.
  $$
\end{lemma}
\begin{proof}
\uses{vertIntBound,limitOfConstant,zeroTendstoDiff,tendsto_zero_Lower,isIntegrable,Perron.tendsto_zero_Upper,RectangleIntegral_tendsTo_VerticalIntegral,tendsto_rpow_atTop_nhds_zero_of_norm_lt_one,integralPosAux,RectangleIntegral,isHolomorphicOn}
\leanok
Let $f(s) = x^s/(s(s+1))$. Then $f$ is holomorphic on the half-plane
  $\{s\in\mathbb{C}:\Re(s)>0\}$. The rectangle integral of $f$ with corners $\sigma-iT$ and
  $\sigma+iT$ is zero. The limit of this rectangle integral as $T\to\infty$ is
  $\int_{(\sigma')}-\int_{(\sigma)}$. Therefore, $\int_{(\sigma')}=\int_{(\sigma)}$.

But we also have the bound $\int_{(\sigma')} \leq x^{\sigma'} * C$, where
  $C=\int_\R\frac{1}{|(1+t)(1+t+1)|}dt$.

Therefore $\int_{(\sigma')}\to 0$ as $\sigma'\to\infty$.

So pulling contours gives $\int_{(\sigma)}=0$.
\end{proof}
